\section{Evaluation}\label{sec:evaluation}

In this section, we revisit the goals of Nugget and evaluate how well we achieve them.

\begin{enumerate}
    \item Analyzing program execution quickly and without restrictions on architecture or host machine.
    \item Automatically generating samples that are independent of the binary.
    \item Quickly validating the selected samples for the targeted benchmarks and input size.
\end{enumerate}

We begin by evaluating the first goal through measurements of the analysis overhead introduced by Nugget, comparing it to both uninstrumented execution and to functional simulation (Section~\ref{sec:interval_analysis_efficiency}). 
We then address the second and third goals by employing two simple sample selection methodologies, random sampling and LLVM IR basic block (IRBB) vector-based k-means sampling, and evaluating their effectiveness (Section~\ref{sec:sampling_methodology_validation}).

\begin{table}[!t]
  \centering

  % ---- Table 1 ----
  \begin{minipage}{\columnwidth}
    \centering
    \scriptsize
    \setlength{\tabcolsep}{4pt}
    \renewcommand{\arraystretch}{1.2}

    \begin{tabular}{|c|c|c|c|}
      \hline
      \textbf{Model} & \textbf{Core} & \textbf{Cache Hierarchy} & \textbf{Memory} \\ \hline
      \makecell[l]{AMD Ryzen\\Threadripper 3960X}
        & \makecell[l]{24-core OOO\\@ 3.8~GHz}
        & \makecell[l]{L1d:32~KiB L1i:32~KiB\\L2:512~KiB L3:128~MiB}
        & \makecell[l]{DDR4 2133~MHz\\136.5~GiB/s} \\ \hline
      \makecell[l]{Ampere Altra}
        & \makecell[l]{160-core OOO\\@ 2.8~GHz}
        & \makecell[l]{L1d:64~KiB L1i:64~KiB\\L2:1~MiB SLC:64~MiB}
        & \makecell[l]{DDR4 3200~MHz\\409.6~GiB/s} \\ \hline
    \end{tabular}

    \caption{Real Hardware Configuration}
    \label{tab:system-configuration}
  \end{minipage}

  \vspace{0.8em}

  % ---- Table 2 ----
  \begin{minipage}{\columnwidth}
    \centering
    \scriptsize

    \begin{tabular}{|p{1cm}|p{2.5cm}|p{2.5cm}|p{1.1cm}|}
      \hline
      \textbf{Suite} & \textbf{Workload} & \textbf{Parallelism} & \textbf{Lang.} \\ \hline
      \makecell[l]{SPEC \\ CPU2017}
        & \makecell[l]{Desktop and server \\workloads}
        & Single-threaded
        & \makecell[l]{C, C++, \\ Fortran} \\ \hline
      \makecell[l]{NPB}
        & HPC kernels
        & Multi-threaded (OMP)
        & Fortran, C \\ \hline
      LSMS
        & \makecell[l]{First-principles \\ materials simulation}
        & Single-threaded
        & C++ \\ \hline
    \end{tabular}

    \caption{Overview of Benchmark Suites Used in Evaluation}
    \label{tab:benchmark-suites}
    
  \end{minipage}

\end{table}
\FloatBarrier % optional, keeps them from drifting past this point

Table~\ref{tab:system-configuration} presents the detailed configurations of the real hardware used in our evaluation. 
For the simulation models, we employ a basic generic functional model using the ATOMIC CPU in gem5, running a guest Ubuntu~24.04 environment, to perform interval analysis on both x86 and Arm architectures.

\revision{
We evaluate Nugget using the SPEC CPU~2017 benchmark suite, the NAS Parallel Benchmarks (NPB) suite~\cite{npb}, and the LSMS benchmark~\cite{LSMS}, covering a wide range of realistic workloads as summarized in Table~\ref{tab:benchmark-suites}. 
}

\revision{We use the ``-O2'' optimization level for the compilation.}
The input chosen for each benchmark is too large to be simulated in feasible time (except for NPB class ``A'').

For SPEC~CPU2017, we use the reference input sizes, which require between 198~billion and 9{,}897~billion instructions.
For NPB, we use class ``C'' and ``D'' inputs, with class ``C'' requiring between 24~billion and 2{,}714~billion instructions and class ``D'' ranging from 1{,}217~billion to 37{,}533~billion instructions. 
For LSMS, we use the Fe input, which requires between 28{,}484~billion and 198{,}885~billion instructions (one binary per ISA).


\shepherd{We used SPEC~CPU2017 Integer suite, except for ``619.lbm\_s'' and ``644.nab\_s'', which are from the Floating-Point suite.}
\revision{
  We do not include CPU~2017 benchmarks that failed to compile with the original SPEC~CPU2017 GCC build system (6 benchmarks) or encountered Fortran linking issues (5 benchmarks).
  We did not encounter the same Fortran linking issues with any NPB Fortran benchmarks.
% All the SPEC CPU~2017 workloads used here, expect for ``619.lbm\_s'' and ``644.nab\_s'', are from the Integer suite.
% The benchmarks not included from the suite either failed to compile with the original SPEC CPU~2017 GCC build system (6 benchmarks) or encountered Fortran linking issues with the LLVM workflow we used (5 benchmarks) while not having issues with NPB's Fortran benchmarks.
% The benchmarks that only has ``x86'' but not ``Arm'' result are crashed due to LLVM backend produces ``x86'' datalayout on ``AArch64''.
% This is a limitation of the LLVM version we used on cross-compiling from an existing LLVM IR file.
}

To accurately measure hardware performance, we employ \textit{cpuset} to strictly bind workloads to specific cores and memory nodes. 
We also disable CPU frequency scaling and fix the CPU at a constant frequency. 
The LLVM passes are implemented on top of LLVM 19.0.0git.
